\scanpage{6}
\mspage{-0.13-}
image of $\varphi'$. Since this image is dense, there are no such $y$'s.
Therefore $\sigma[\varphi^{-1}(U)]=\varphi^{-1}(U)\times_U U'$, and $\varphi'$ is surjective.
Therefore $\overline\omega$ is geometrically injective without restriction, and by \uline{Zariski's Main Theorem}, $\overline\omega$ is an isomorphism. \uline{QED}
\margin{!}

\uline{Remark} \uline{(5)}. The following are equivalent:

(a) $(Y,\varphi)$ is a universal categorical quotient of $X$ by $G$,

(b) for all affine schemes $Y'$, and morphisms $Y'\longrightarrow Y$, if $\varphi':X'\to Y'$ is the base extension, then $(Y',\varphi')$ is a categorical quotient of $X'$ by $G$,

(c) there is an affine open covering $\{U_i\}$ of $Y$ such that if $\varphi_i:\varphi^{-1}(U_i)\longrightarrow U_i$ is the restriction of $\varphi$, then $(U_i,\varphi_i)$ is a universal categorical quotient of $\varphi^{-1}(U_i)$ by $G$, for all $i$.

(6). In Chapter 1, we will need a criterion for $(Y,\varphi)$ to be a categorical quotient which does not imply that $(Y,\varphi)$ is a geometric quotient. In fact, the following 3 conditions suffice:

\[
\text{i)}\qquad\varphi\circ\sigma=\varphi\circ p_2,
\]
ii) $\underline o_Y$ is the subsheaf of invariants of $\varphi_*(\underline o_X)$,

iii) if $W$ is an invariant closed subset of $X$, then $\varphi(W)$ is closed in $Y$; if $W_i$, $i\in I$, is a set of invariant closed subsets of $X$, then:
\[
\varphi\left(\bigcap_{i\in I}W_i\right)=\bigcap_{i\in I}\varphi(W_i)
\]
\margin{\diagram{\begin{tikzcd}[column sep=large]
R\arrow[d,shift right=.6ex,"p_1"']\arrow[d,shift left=.6ex,"p_2"] & \\
\varphi:X\arrow[r,"\unsure{f}"'] & Y
\end{tikzcd}}
morphism of preschemes \unsure{st.h.} $\varphi p_1=\varphi p_2$ \unsure{with} $Y$ is a quotient $X/R$ if
(i) $\underline O_Y$ is the sheaf of invariant of
\diagram{\begin{tikzcd}[column sep=small]
\underline O_Y\arrow[r,dash] & f_*(\underline O_X)\arrow[r,shift left=.6ex]\arrow[r,shift right=.6ex] & g_*(\unsure{O_R})
\end{tikzcd}}
(ii) Any invariant open $U$ of $X$ is the inverse image of an open subset of $Y$.}
\margin{\figure{\begin{tikzpicture}\draw[thin] (0,0) .. controls (.09,-.6) and (.07,-1.05) .. (-.3,-1.65);\end{tikzpicture}}Il paraît suffisant que $\varphi$ \unsure{soit} \uline{\unsure{quasi-qct|quasi-cpt}} par \unsure{exp.}}
\mspage{5}
